\section{Discussion}
\label{sec:discussion}

\subsection{From optimization pipelines to adaptive scientific research}

A fixed optimizer solves a predefined problem. A research agent must additionally decide whether the current problem is the right one. In tendon-driven robot design, the limiting mechanism may change as the structure evolves: delayed tendon recruitment may be replaced by insufficient high-load stiffness, dynamic settling, or a routing-induced loss of force authority. A system that only repeats the same optimization cannot necessarily respond to such changes.

The Fixed-Plan comparison is therefore central to C1. It tests whether evidence-driven changes to the research route add value beyond access to the same mathematical tools. The One-shot LLM Plan condition, if added, further distinguishes initial language-model planning from repeated evidence-driven replanning.

\subsection{Why research-intent projection matters}

The central methodological idea in C2 is not merely exposing more numerical tools. It is the projection from a high-level physical question into a concrete mathematical problem. The agent should be able to state what mechanism it suspects, which parameters may influence it, and which physical observables would reveal the effect. The fixed program should then determine whether those quantities are registered, whether the proposed inner problem is well-defined, which mathematical method is applicable, and what numerical evidence can be returned.

This separation avoids two extremes. In one, the LLM guesses structures and parameters without quantitative physical support. In the other, the LLM becomes a thin wrapper around a fixed optimizer and cannot change the scientific question. Research-intent projection provides a middle ground in which the agent chooses the question while deterministic mathematics owns the numerical semantics.

\subsection{Mechanism sensitivity and task capability are complementary}

A mechanism-level objective can help discover a promising structural relationship, but it may not improve the task. Conversely, a task-level capability measure can show whether a design change improves the best achievable task outcome, but may not explain why. The two should be used together: mechanism sensitivity proposes a physically interpretable direction, while task-conditioned capability sensitivity tests whether that direction improves the task after control re-optimization.

The distinction is important for passive structures as well. Differentiating tendon geometry at a fixed configuration is not the same as differentiating a re-equilibrated passive response. The latter requires an equilibrium sensitivity calculation and careful treatment of unilateral active-set changes.

\subsection{Evidence as part of the research method}

Knowledge and evidence organization are not merely logging utilities. They influence which observations are available to the agent, whether earlier findings can be reused, and whether a later decision can be justified. Reporter interpretation is useful only if it remains faithful to fixed numerical evidence and clearly separates observation from explanation. The No-Knowledge and Reporter ON/OFF experiments therefore evaluate distinct aspects of C3.

The effective research action ratio complements physical task metrics by measuring whether research actions produce usable scientific evidence. It must be assessed with a frozen rubric; otherwise, it risks rewarding persuasive descriptions rather than genuine progress.

\subsection{Limitations and future work}

Several questions remain open:
\begin{itemize}
  \item \textbf{Intent-projection implementation:} the exact boundary of the current \texttt{scientific\_question\_v4} pathway must be audited before the No-Intent-Projection ablation can be treated as a clean component intervention.
  \item \textbf{Sensitivity guarantees:} the envelope-theorem relation requires local optimality and regularity. The current workflow does not provide a universal gradient-error bound or global descent certificate.
  \item \textbf{Active-set transitions:} slack-to-taut switching and contact may produce non-smooth response maps. Directional and mode-specific analyses are needed in these regimes.
  \item \textbf{Knowledge alignment:} the historical P3 knowledge assets and the current Core knowledge snapshot must be aligned before a strict No-Knowledge comparison.
  \item \textbf{Probabilistic diagnosis:} the system records expected and weakening observations but does not yet implement calibrated Bayesian information gain or numerical value-of-information optimization.
  \item \textbf{Repeated runs:} single-seed pilots are exploratory. Quantitative comparisons require independent repetitions and uncertainty estimates.
  \item \textbf{Hardware validation:} the current evidence is primarily simulation-based. Physical prototypes are needed to study tendon friction, manufacturing tolerances, actuator dynamics, and sim-to-real effects.
\end{itemize}

The immediate priority is to evaluate adaptive research decisions under matched budgets, then isolate mathematical and evidence capabilities. Only after these comparisons should the system be described as improving research efficiency or generalizing across LLM backbones.
